% ===========================================================================
% SPRINGER LNCS CONFERENCE PAPER — AQUAPONICAI
% Authors: Chandan K, Dr. Deepika K | Dept. MCA, RVCE, Bengaluru
%
% OVERLEAF INSTRUCTIONS:
%  1. New Project → Blank Project → paste this file as main.tex
%  2. Compiler: pdfLaTeX
%  3. Upload 4 images to Overleaf root before compiling:
%       fig_arch.png      (System architecture diagram)
%       fig_survey.png    (Screenshot: voice survey recording state)
%       fig_scenarios.png (Screenshot: Scenario Analysis bar chart)
%       fig_dashboard.png (Screenshot: main dashboard overview)
% ===========================================================================

\documentclass[runningheads]{llncs}

\usepackage[T1]{fontenc}
\usepackage{amsmath}
\usepackage{amssymb}
\usepackage{graphicx}
\usepackage{booktabs}
\usepackage{array}
\usepackage{ragged2e}
\usepackage{url}
% hyperref must be loaded last
\usepackage[hidelinks]{hyperref}

% Safe INR notation
\newcommand{\inr}[1]{Rs.\,#1}

% ── Title ───────────────────────────────────────────────────────────────────
\title{AquaponicAI: A Voice-Driven Adaptive Decision Support System
for Aquaponic and Land Farm Planning with Multi-Scenario
Financial Analysis}

\titlerunning{AquaponicAI: Voice-Driven Decision Support for Farm Planning}

% ── Authors ─────────────────────────────────────────────────────────────────
\author{Chandan K\inst{1} \and Dr.~Deepika K\inst{1}}

\authorrunning{C.~K and D.~K}

\institute{
  Department of MCA, RV College of Engineering,\\
  Bengaluru, Karnataka, India\\
  \email{\{chandank.mca24, deepikak\}@rvce.edu.in}
}

\begin{document}
\maketitle

% ── Abstract ────────────────────────────────────────────────────────────────
\begin{abstract}
Aquaponic and land-based farm operators in developing economies face a
persistent challenge: translating operational knowledge into structured
financial plans without access to domain experts or expensive software.
This paper presents AquaponicAI, an end-to-end decision-support platform
that enables farm operators to generate multi-scenario financial projections
through a voice-first adaptive survey interface. The system provides two
survey modules: an aquaponics-specific questionnaire covering fish, crop,
water quality, and financial parameters, and a land-and-crop planning
module with real-time market price integration via the Government of India
Agmarknet API. Voice input is processed through a dual-path transcription
pipeline combining browser-native Web Speech API with server-side
\texttt{faster-whisper}, followed by domain vocabulary normalization and
confidence-aware entity extraction. A stateless questionnaire engine with
conditional skip logic drives structured data collection across six question
types. Collected data feeds a deterministic financial planning engine that
computes capital expenditure, operating costs, return on investment, net
present value, internal rate of return, payback period, break-even month,
and per-area productivity across three revenue scenarios. The platform is
implemented using FastAPI on Python~3.12, React~18, PostgreSQL~16, and
Redis~7, fully containerized via Docker Compose. A 36-case validation suite
(30 financial and questionnaire-engine unit tests, 6 voice-interpretation
unit tests) was executed with full pass rate in our evaluation run.

\keywords{Aquaponics \and Decision support system \and Voice user interface
\and Adaptive questionnaire \and Multi-scenario financial analysis
\and Speech recognition \and Smart agriculture}
\end{abstract}


% ── I. INTRODUCTION ─────────────────────────────────────────────────────────
\section{Introduction}
\label{sec:intro}

Aquaponics, a recirculating system integrating aquaculture with
hydroponics~\cite{ref1,ref21,ref22}, produces fish and crops in a shared
water circuit, achieving water savings of 90--95\% compared to
conventional soil farming while generating dual revenue
streams~\cite{ref7}. Despite
strong agronomic advantages, smallholder aquaponic operations face high
failure rates driven principally by inadequate financial planning.
Operators typically rely on generic spreadsheet templates that do not
model the co-dependent revenue cycles of fish and crop production, nor
communicate how returns change under adverse market or yield conditions.

Land-based smallholder farming faces analogous challenges. Crop
selection, cost structure, and price volatility interact in ways that
are difficult to anticipate without quantitative scenario modelling.
Studies of smallholder agri-enterprises identify lack of structured
financial planning tools as a primary contributor to operational
losses within the first three years~\cite{ref4,ref35}.
In India, where smallholder farming supports hundreds of millions of
livelihoods, digital planning tools have seen limited penetration
beyond crop advisory services~\cite{ref18}.

Decision support systems (DSS) for precision agriculture have been
researched extensively, with applications spanning crop disease
identification, yield prediction, and irrigation
scheduling~\cite{ref2,ref15}. However, adoption among aquaponic and
smallholder land operators remains low due to two structural barriers.
First, existing tools require form-based interfaces that impose
cognitive burden on operators managing physically demanding daily
operations. Second, financial projections are typically single-scenario
estimates that convey no information about downside risk, limiting
their utility for investment decisions.

Voice-first interfaces present a compelling solution to the interaction
barrier. Studies on agricultural advisory systems demonstrate that
speech-based data entry reduces task completion time and cognitive
load compared to form-based input, particularly for operators with
limited digital literacy~\cite{ref3,ref20,ref30}. However, agricultural speech
processing introduces domain-specific challenges---species name
homophones, spoken number conventions, and Indian English accent
variation---that general-purpose ASR engines do not address.

This paper presents \textbf{AquaponicAI}, a unified platform that
addresses both barriers. The primary contributions are:

\begin{enumerate}
  \item A dual-module adaptive questionnaire platform with voice
        input, six question types, conditional skip logic, and
        real-time market price integration for crop revenue estimation.

  \item A domain-specific voice interpretation pipeline with
        aquaponics vocabulary normalization, spoken-number normalisation
        for Indian large-number vocabulary, and a confidence-fusion
        confirmation gate designed to reduce unnecessary user interruptions.

  \item A deterministic multi-scenario financial planning engine
        computing 14 financial metrics across pessimistic
        (revenue $\times$0.75), base, and optimistic
        (revenue $\times$1.30) scenarios, with 30 unit tests
        verifying correctness of all metric computations and
        live-recalculation for assumption sensitivity exploration.

  \item A containerized full-stack deployment (Docker Compose,
        Nginx reverse proxy) designed to minimise infrastructure
        requirements and tolerate intermittent connectivity, with
        the full source available for replication.
\end{enumerate}


% ── II. RELATED WORK ────────────────────────────────────────────────────────
\section{Related Work}
\label{sec:related}

\subsection{AI Decision Support in Precision Agriculture}

Machine learning and AI-based DSS platforms for smart agriculture have
been reviewed extensively~\cite{ref2,ref15,ref23,ref24,ref27,ref32}.
Applications range from crop disease detection to supply chain
optimization. Big data analytics and deep learning are increasingly
central to farm management~\cite{ref23,ref24}, while Agriculture~5.0
frameworks emphasise human-centred intelligent platforms~\cite{ref32}.
Despite this breadth, platforms designed for the financial planning
needs of smallholder operators remain scarce, and those that exist
typically target crop-only farming rather than the dual revenue model
of aquaponics.

\subsection{Voice Interfaces for Agricultural Advisory}

Voice-based advisory systems for low-resource agricultural contexts
have demonstrated feasibility across South Asia~\cite{ref3,ref20}. Key
challenges include ASR robustness to regional accents, domain-specific
vocabulary, and graceful recovery from low-confidence transcriptions.
Named entity recognition (NER) for agricultural knowledge extraction,
supported by pre-trained bidirectional language models~\cite{ref25},
has been explored to ground free-form speech input to structured farm
parameters~\cite{ref9}. The transformer architecture~\cite{ref29}
underpins modern self-supervised speech models such as
wav2vec\,2.0~\cite{ref10} and Whisper~\cite{ref5}, which deliver
strong general transcription accuracy. Advances in speech data
augmentation~\cite{ref28} and deep recurrent sequence
modelling~\cite{ref26} have further improved robustness, yet
domain-specific post-processing remains necessary to handle
out-of-vocabulary terminology and regional accent variation. Existing agricultural advisory systems rely on cloud
ASR without such a layer; AquaponicAI contributes a structured
vocabulary replacement pipeline and dual-threshold confirmation logic
that reduces unnecessary confirmation dialogs.

\subsection{Financial Planning Tools for Farm Enterprises}

Multi-scenario financial modelling using NPV and IRR is standard
in capital budgeting for agri-enterprises~\cite{ref4,ref33}. Economic
analyses of aquaponic ventures specifically highlight the importance
of sensitivity modelling across revenue scenarios~\cite{ref31}. Machine
learning approaches to commodity price forecasting improve revenue
projection accuracy in agri-financial models~\cite{ref11}. Market
price datasets from Indian agricultural markets have been used to
benchmark crop price prediction systems~\cite{ref17}. Digital tools
that automate this process for smallholder operators are emerging in
agri-fintech but typically target crop-only farming. To the best of
our knowledge, the integration of voice-driven data collection,
adaptive survey methodology, real-time market price lookup, and
scenario-based financial analysis in a single pipeline for both
aquaponic and land farming has not been reported in the surveyed
literature.

\subsection{Speech Recognition for Domain Applications}

Large-scale transformer-based ASR models such as
Whisper~\cite{ref5} have demonstrated strong general-purpose
transcription accuracy, building on foundational work in deep
recurrent speech modelling~\cite{ref26} and transformer
architectures~\cite{ref29}. Data augmentation strategies such as
SpecAugment~\cite{ref28} have further improved robustness to
acoustic variation. Domain-specific vocabulary and Indian-accented
English remain challenging for out-of-the-box models~\cite{ref10}.
AquaponicAI addresses this with a structured post-processing layer
rather than model fine-tuning, enabling deployment without
GPU-intensive training.


% ── III. SYSTEM ARCHITECTURE ────────────────────────────────────────────────
\section{System Architecture}
\label{sec:arch}

\begin{figure}[tbp]
  \centering
  \includegraphics[width=\textwidth]{fig_arch}
  \caption{AquaponicAI three-tier architecture: React\,18 presentation
    layer, FastAPI application layer, and PostgreSQL\,16 + Redis data
    layer, orchestrated via Docker Compose behind an Nginx reverse proxy.}
  \label{fig:arch}
\end{figure}

AquaponicAI follows a three-tier architecture (Fig.~\ref{fig:arch}).

\paragraph{Presentation Layer.}
A React\,18 SPA bundled with Vite serves two survey modules
(Aquaponics Survey and Land Survey), a Financial Plan Dashboard,
Farm Records management, and an Analytics view. Zustand manages
global state including session context and financial analysis results.
A custom \texttt{useVoiceRecorder} hook wraps the browser-native
Web Speech API, providing real-time interim transcript events.
``AI~Voice~On'' and ``Validation~On'' toggles allow users to enable
or disable voice readback and server-side answer validation per
session. An Axios HTTP client with JWT interceptors handles all
API communication.

\paragraph{Application Layer.}
A FastAPI application on Python\,3.12 uses fully asynchronous I/O
via \texttt{asyncio} and SQLAlchemy\,2.x. Thin routers delegate
immediately to service modules: (i)~session and questionnaire
management, (ii)~voice transcription via
\texttt{faster-whisper}~\cite{ref5}, (iii)~financial planning and
scenario computation, (iv)~real-time market price lookup via the
Government of India Agmarknet API with fallback to benchmark prices,
and (v)~Google Sheets export.

\paragraph{Data Layer.}
PostgreSQL\,16 uses UUID primary keys throughout. \texttt{JSONB}
columns with GIN indexes store survey context and financial plan
scenarios, enabling schema-flexible answer accumulation without
migrations. Redis\,7 provides response caching for market price
lookups.

All services are orchestrated via Docker Compose with Nginx routing
\texttt{/api} to the backend and all other traffic to the React SPA.

\begin{figure}[tbp]
  \centering
  \includegraphics[width=0.92\textwidth]{fig_dashboard}
  \caption{AquaponicAI main dashboard providing access to Aquaponics
    Survey, Land Survey, Farm Records, Reports, and the Farm Profit
    Trend chart.}
  \label{fig:dashboard}
\end{figure}


% ── IV. METHODOLOGY ─────────────────────────────────────────────────────────
\section{Methodology}
\label{sec:method}

\subsection{Adaptive Questionnaire Engine}

The \texttt{QuestionnaireEngine} manages a bank of 29 questions
organized across six operational categories (Table~\ref{tab:qbank}).
The engine is stateless: all session progress accumulates in a
\texttt{context\_data} JSONB column passed into each call to
\texttt{get\_next\_question()}.

\begin{table}[tbp]
  \caption{Questionnaire Bank: Categories and Coverage}
  \label{tab:qbank}
  \centering
  \small
  \setlength{\tabcolsep}{5pt}
  \begin{tabular}{llc}
    \toprule
    \textbf{Category} & \textbf{Key Parameters Captured} & \textbf{Qs} \\
    \midrule
    Farm Setup     & Name, location, area, system type          & 5  \\
    Fish Data      & Species, count, tank vol., feed, cycle     & 6  \\
    Crop Data      & Crop types, grow area, monthly yield       & 3  \\
    Water Quality  & pH, dissolved oxygen, temperature          & 3  \\
    Financial      & CAPEX, OPEX items, revenue, horizon        & 10 \\
    Goals          & Primary goal, operational challenge        & 2  \\
    \midrule
    \textbf{Total} & \textbf{6 categories, 6 question types}    & \textbf{29} \\
    \bottomrule
  \end{tabular}
\end{table}

The six question types are TEXT, NUMBER, SELECT, MULTISELECT,
BOOLEAN, and DATE. Questions carry an optional
\texttt{follow\_up\_condition} field; for example, the crop-yield
question is skipped when monthly revenue is directly provided.
A regex-based condition evaluator avoids code execution during
skip logic evaluation.

Type-aware parsing in \texttt{parse\_answer()} handles:

\paragraph{NUMBER.}
Digit-first extraction with fallback to a word-to-digit parser
covering Indian large-number vocabulary
(\emph{lakh}\,$\rightarrow$\,100\,000;
\emph{crore}\,$\rightarrow$\,10\,000\,000)
and common STT homophones
(\emph{won}\,$\rightarrow$\,1; \emph{ate}\,$\rightarrow$\,8;
\emph{to/too}\,$\rightarrow$\,2).

\paragraph{BOOLEAN.}
Positive/negative vocabulary matching.

\paragraph{SELECT / MULTISELECT.}
Fuzzy substring matching with free-form fallback.

\paragraph{TEXT (location).}
India city/state validation against a 36-item State/UT vocabulary
with spoken comma-less input handling.

\subsection{Voice Interpretation Pipeline}

\begin{figure}[tbp]
  \centering
  \includegraphics[width=0.92\textwidth]{fig_survey}
  \caption{Voice survey in active recording state with the ``AI~Voice~On''
    and ``Validation~On'' toggles and the captured answer shown.}
  \label{fig:survey}
\end{figure}

Voice input proceeds through three sequential stages after capture
(Fig.~\ref{fig:survey} shows the user-facing interface).

\paragraph{Stage 1: Filler Removal and Domain Normalization.}
\texttt{post\_process\_transcript()} removes 12 filler tokens
(\emph{um, uh, erm, like, you know, i mean, actually, basically,
sort of, kind of, right, literally}) via compiled regex, then applies
domain replacement rules. Aquaponics-specific substitutions include
species corrections (\emph{till, talapia, to lapia} $\rightarrow$
tilapia; \emph{troat, trowt} $\rightarrow$ trout;
\emph{barry, mundi, baramundi} $\rightarrow$ barramundi) and system
type corrections (\emph{nft, an empty, n f t} $\rightarrow$ nft;
\emph{media bead/bid/bad} $\rightarrow$ media bed). This domain
normalization layer addresses the well-documented difficulty of
general-purpose ASR models with out-of-vocabulary domain
terms~\cite{ref5,ref10}.

\paragraph{Stage 2: Spoken Number Normalisation and Answer Confidence Scoring.}
Agricultural voice input in India routinely uses large-scale number
vocabulary not present in standard ASR language models.
\texttt{parse\_answer()} applies a four-stage normalisation cascade
(Table~\ref{tab:numstages}) to any NUMBER-type response before
validation~\cite{ref9,ref10}.

\begin{table}[tbp]
  \caption{Spoken Number Normalisation Priority Chain}
  \label{tab:numstages}
  \centering
  \small
  \setlength{\tabcolsep}{5pt}
  \begin{tabular}{clp{5.5cm}}
    \toprule
    \textbf{Stage} & \textbf{Method} & \textbf{Examples handled} \\
    \midrule
    1 & Digit regex          & \texttt{"5000"}, \texttt{"5,000"} \\
    2 & Indian scale words   & \emph{lakh}$\!\to\!$1\,e5,
                               \emph{crore}$\!\to\!$1\,e7 \\
    3 & STT homophones       & \emph{won}$\!\to\!$1,
                               \emph{ate}$\!\to\!$8,
                               \emph{to/too}$\!\to\!$2,
                               \emph{for}$\!\to\!$4 \\
    4 & Ordinal words        & \emph{twenty thousand}$\!\to\!$20\,000 \\
    \bottomrule
  \end{tabular}
\end{table}

The normalised value is $v = d \cdot m_s$, where $d$ is the digit
part extracted by Stage~1--4 and $m_s \in \{1,\,10^3,\,10^5,\,10^7\}$
is the scale multiplier from Indian number vocabulary
(\emph{thousand}, \emph{lakh}, \emph{crore} respectively).
Parsing success probability serves as answer-level confidence:
\begin{equation}
  \hat{e} = \begin{cases}
    1.0  & \text{Stage 1 (digit) match} \\
    0.85 & \text{Stage 2 (scale word) match} \\
    0.70 & \text{Stage 3--4 (homophone/ordinal) match} \\
    0.30 & \text{no match (free-text fallback)}
  \end{cases}
  \label{eq:econf}
\end{equation}
For non-NUMBER types, $\hat{e}$ is $1.0$ on exact option match,
$0.75$ on fuzzy substring match, and $0.40$ on free-text fallback.
This score feeds into the confidence fusion gate.

\paragraph{Stage 3: Confidence Fusion and Confirmation Gate.}
STT confidence $c_{\mathrm{stt}} \in [0,1]$ is fused with entity
confidence $\hat{e}$:
\begin{equation}
  c_{\mathrm{comb}} =
    \mathrm{clamp}_{[0,1]}\!\left(0.6\,c_{\mathrm{stt}} + 0.4\,\hat{e}\right)
  \label{eq:fusion}
\end{equation}
The 0.6/0.4 weighting reflects that $c_{\mathrm{stt}}$ is a
calibrated probability estimate from a neural model~\cite{ref10},
whereas $\hat{e}$ is a coarser rule-based signal.
A confirmation dialog is triggered when:
\begin{equation}
  \mathrm{confirm} =
    (c_{\mathrm{comb}} < \theta_{e}) \;\lor\;
    (c_{\mathrm{stt}} < \theta_{\mathrm{stt}})
  \label{eq:confirm}
\end{equation}
with default thresholds $\theta_{e}=0.65$ and
$\theta_{\mathrm{stt}}=0.45$, selected heuristically to balance
false-confirmation rate against missed-error rate; both are
configurable via environment variables and subject to empirical
calibration in the planned field-trial phase.

The server-side transcription path uses \texttt{faster-whisper}
(default model: \texttt{large-v3}) with lazy initialization and
optional GPU acceleration. The browser-side path uses the Web Speech
API for zero-upload-latency real-time streaming. Both paths feed into
the same interpretation pipeline.

\subsection{Financial Planning Engine}

Survey responses are mapped to a \texttt{FinancialInputs} dataclass
comprising three CAPEX components (infrastructure, equipment, initial
stock), five monthly OPEX components (feed, labour, utilities,
maintenance, other), and three revenue streams (fish/crops/other for
aquaponics; per-crop revenue derived from yield and market price for
land farming).

\paragraph{Core Financial Metrics.}
\begin{align}
  \mathrm{CAPEX} &= C_{\mathrm{infra}} + C_{\mathrm{equip}}
                    + C_{\mathrm{stock}} \label{eq:capex}\\
  \mathrm{ROI}   &= \frac{12(f_s R_m - \mathrm{OPEX}_m)}
                          {\mathrm{CAPEX}} \times 100\%
                    \label{eq:roi}
\end{align}
where $f_s \in \{0.75, 1.00, 1.30\}$ is the scenario revenue factor
for pessimistic ($-25\%$), base, and optimistic ($+30\%$) scenarios
respectively. These bounds are consistent with the downside/upside
sensitivity ranges reported in aquaponic financial risk
modelling~\cite{ref4}, where asymmetric ranges reflect that revenue
upside from dual fish--crop streams exceeds typical downside from
single-product exposure.

\paragraph{Net Present Value.}
\begin{equation}
  \mathrm{NPV} = -\mathrm{CAPEX} +
    \sum_{m=1}^{H}
      \frac{f_s R_m - \mathrm{OPEX}_m}{(1+r_m)^{m}}
  \label{eq:npv}
\end{equation}
where $H$ is the planning horizon in months,
$r_m = (1+r_a)^{1/12}-1$ is the monthly discount rate,
and $r_a = 0.08$ (8\% annual).

\paragraph{Internal Rate of Return.}
IRR is approximated by bisection search (maximum 200 iterations,
convergence tolerance $|\mathrm{NPV}(r)| < 10^{-4}$) for $r^*$
satisfying $\mathrm{NPV}(r^*)=0$, then annualised:
\begin{equation}
  \mathrm{IRR}_a = (1 + r^*_m)^{12} - 1
  \label{eq:irr}
\end{equation}

\paragraph{Payback Period.}
\begin{equation}
  T_{\mathrm{pay}} = \frac{\mathrm{CAPEX}}{f_s R_m - \mathrm{OPEX}_m},
  \quad f_s R_m > \mathrm{OPEX}_m
  \label{eq:payback}
\end{equation}
Break-even month $B$ is the first $m$ at which cumulative cash flow
crosses zero.

\paragraph{Market Price Integration.}
For the Land Survey module, monthly crop revenue is computed as
yield\,$\times$\,price/kg, where price is fetched in real time from
the Agmarknet commodity price API (data.gov.in). The system normalizes
district and state names (``Bangalore'' $\rightarrow$ ``Bengaluru
Urban'', ``KA'' $\rightarrow$ ``Karnataka''), retries with
progressively relaxed geographic filters on empty results, and falls
back to benchmark prices on API failure. This enables the financial
plan to reflect prevailing market conditions rather than static
assumptions~\cite{ref11,ref17}.

\paragraph{Recommendation Engine.}
A deterministic rule-based engine evaluates base-scenario metrics
against eight financial health rules, generating structured
recommendations with category, priority, and actionable detail
fields, drawing on principles from recommendation systems
research~\cite{ref34}.
Rules cover: cash-flow sign, operating margin (10\% threshold), feed
cost share (45\%), labour share (45\%), utility share (25\%), payback
horizon (48 months), revenue concentration risk (85\% single stream),
and NPV sign at horizon end.

\paragraph{Live Recalculation.}
The frontend provides an ``Adjust Assumptions'' panel with sliders
for key revenue and cost parameters. Changes trigger synchronous
recalculation and immediate re-rendering of all metrics and charts
without a server request, enabling rapid sensitivity exploration.


% ── V. SYSTEM VALIDATION ────────────────────────────────────────────────────
\section{System Validation}
\label{sec:eval}

Validation follows a software-correctness methodology suited to a
deterministic computational system prior to live user trials. The test
suite verifies that the system produces expected outputs under all
defined fixture conditions; real-world accuracy under field inputs
remains a subject for planned user trials.

\subsection{Test Suite Composition}

\noindent\textbf{Financial unit tests}
(\texttt{test\_financial\_service.py}): 30 functions covering CAPEX
and OPEX computation, ROI, payback, break-even, NPV, and IRR
calculation; scenario monotonicity across all three factors; per-area
metric derivation; revenue and cost share percentage sum-to-100
assertions; zero-revenue edge case; and infinite-payback detection.

\smallskip
\noindent\textbf{Voice interpretation tests}
(\texttt{test\_voice\_interpretation.py}): 6 functions covering
spoken-number normalisation across all four cascade stages
(Table~\ref{tab:numstages}), Indian scale-word parsing (\emph{lakh},
\emph{crore}), STT homophone correction, filler removal,
threshold-triggered confirmation, and confidence fusion arithmetic
verified against \eqref{eq:fusion}.

\smallskip
\noindent\textbf{STT evaluation harness}
(\texttt{stt\_evaluation\_harness.py}): 4 representative cases
covering common ASR failure modes: prefix conversational noise,
all-lowercase output, filler-laden input, and low-confidence
transcript output. All four cases confirmed correct post-processing
behaviour in our evaluation run.

\subsection{Reference Fixture}

The primary fixture used across all 30 financial tests sets:
\begin{itemize}
  \item Infrastructure: \inr{300{,}000}; Equipment: \inr{150{,}000};
        Initial stock: \inr{50{,}000}
        $\Rightarrow$ CAPEX\,=\,\inr{500{,}000}
  \item Monthly OPEX: \inr{30{,}000}
        (feed \inr{8k}\,+\,labour \inr{15k}\,+\,utilities
         \inr{5k}\,+\,maintenance \inr{2k})
  \item Monthly revenue: \inr{55{,}000}
        (fish \inr{40k}\,+\,crops \inr{15k})
  \item Farm area: 1{,}000\,m\textsuperscript{2};
        Horizon: 24 months
\end{itemize}

All expected outputs are analytically derivable from these inputs,
permitting exact numeric assertions with \texttt{pytest.approx}.


% ── VI. RESULTS AND ANALYSIS ────────────────────────────────────────────────
\section{Results and Analysis}
\label{sec:results}

\subsection{Financial Engine Correctness}

Table~\ref{tab:fin} reports computed values against analytically
derived expected values for the reference fixture. All 30 financial
tests passed in our evaluation run.

\begin{table}[tbp]
  \caption{Financial Engine Unit Test Results (Reference Fixture)}
  \label{tab:fin}
  \centering
  \small
  \setlength{\tabcolsep}{5pt}
  \begin{tabular}{lccl}
    \toprule
    \textbf{Metric} & \textbf{Expected} & \textbf{Computed} & \\
    \midrule
    Total CAPEX                & \inr{5,00,000} & \inr{5,00,000} & \checkmark \\
    Monthly OPEX               & \inr{30,000}   & \inr{30,000}   & \checkmark \\
    Monthly Revenue (base)     & \inr{55,000}   & \inr{55,000}   & \checkmark \\
    Monthly Net CF (base)      & \inr{25,000}   & \inr{25,000}   & \checkmark \\
    Annual Net Profit (base)   & \inr{3,00,000} & \inr{3,00,000} & \checkmark \\
    ROI \eqref{eq:roi}         & 60.0\%         & 60.0\%         & \checkmark \\
    Payback \eqref{eq:payback} & 20\,mo         & 20\,mo         & \checkmark \\
    Break-even month           & 20             & 20             & \checkmark \\
    NPV (base, 24\,mo)         & $>$0           & $>$0           & \checkmark \\
    Scenario order             & P$<$B$<$O      & P$<$B$<$O      & \checkmark \\
    Profit margin              & 45.45\%        & 45.45\%        & \checkmark \\
    OPEX/Revenue               & 54.55\%        & 54.55\%        & \checkmark \\
    Revenue/m$^{2}$/yr         & \inr{660}      & \inr{660}      & \checkmark \\
    Profit/m$^{2}$/yr          & \inr{300}      & \inr{300}      & \checkmark \\
    \bottomrule
    \multicolumn{4}{l}{\footnotesize
      P\,=\,Pessimistic ($f_s{=}0.75$); B\,=\,Base; O\,=\,Optimistic
      ($f_s{=}1.30$).}\\
    \multicolumn{4}{l}{\footnotesize
      ROI\,=\,300k/500k\,$\times$\,100\,=\,60\%.
      Payback\,=\,500k/25k\,=\,20\,months.}\\
  \end{tabular}
\end{table}

The ROI of 60\% is analytically exact: annual net profit =
\inr{25{,}000}\,$\times$\,12\,=\,\inr{3{,}00{,}000}; CAPEX =
\inr{5{,}00{,}000}; ROI = 300k/500k\,$\times$\,100 = 60\%.
Payback = 500k/25k = 20 months. Profit margin = 300k/660k = 45.45\%.
Scenario monotonicity is confirmed: pessimistic ROI = 45\%
(revenue\,$\times$0.75, net = \inr{11{,}250}/mo) and optimistic
ROI = 78\% (revenue\,$\times$1.30, net = \inr{41{,}500}/mo).

\subsection{Voice Interpretation Results}

Table~\ref{tab:voice} reports all six voice interpretation test
outcomes.

\begin{table}[tbp]
  \caption{Voice Interpretation Unit Test Results}
  \label{tab:voice}
  \centering
  \small
  \setlength{\tabcolsep}{5pt}
  \begin{tabular}{
    >{\RaggedRight\arraybackslash}p{3.0cm}
    >{\RaggedRight\arraybackslash}p{3.2cm}
    >{\RaggedRight\arraybackslash}p{3.0cm}
    c}
    \toprule
    \textbf{Test Case} & \textbf{Input} & \textbf{Expected} & \\
    \midrule
    Clean transcript
      & $c_{\mathrm{stt}} > 0.45$
      & Extract; no confirm
      & \checkmark \\[2pt]
    Low STT confidence
      & $c_{\mathrm{stt}} < 0.45$
      & Confirm triggered
      & \checkmark \\[2pt]
    Prefix noise
      & ``My farm is called...''
      & Prefix stripped
      & \checkmark \\[2pt]
    Filler tokens
      & ``um uh Green Ridge''
      & Fillers removed
      & \checkmark \\[2pt]
    Domain vocab
      & ``till'' in input
      & Mapped to tilapia
      & \checkmark \\[2pt]
    Spoken number
      & ``twenty thousand''
      & Parsed to 20000
      & \checkmark \\
    \bottomrule
  \end{tabular}
\end{table}

\subsection{Multi-Scenario Financial Analysis}

Fig.~\ref{fig:scenarios} shows the financial analysis output as
rendered in the platform. The Scenario Analysis chart presents
projected annual profit across six scenarios, giving operators a
complete risk envelope. The donut charts decompose cost structure
and revenue mix, enabling identification of dominant cost and
revenue drivers before investment decisions.

\begin{figure}[tbp]
  \centering
  \includegraphics[width=0.92\textwidth]{fig_scenarios}
  \caption{Scenario Analysis view: horizontal bar chart of projected
    annual profit across six scenarios, monthly cost breakdown donut,
    and revenue mix donut.}
  \label{fig:scenarios}
\end{figure}

\subsection{Market Price Integration}

The Land Survey module fetches live commodity prices from the
Agmarknet API. In a representative evaluation, prices for tomato and
capsicum were retrieved from Bengaluru Urban (Karnataka). The system
normalizes state aliases (``KA'' $\rightarrow$ ``Karnataka'') and
district names (``Bangalore'' $\rightarrow$ ``Bengaluru Urban''),
infers state from known district names, and falls back to benchmark
prices after two failed API attempts. These normalisation and fallback
behaviours are verified by four dedicated unit tests.


% ── VII. DISCUSSION AND CONCLUSION ──────────────────────────────────────────
\section{Discussion and Conclusion}
\label{sec:conclusion}

\subsection{Discussion}

The AquaponicAI prototype indicates that voice-driven data collection,
adaptive survey logic, real-time market price integration, and
deterministic multi-scenario financial analysis can be unified in a
single deployable pipeline at modest infrastructure cost.

\paragraph{Stateless engine architecture.}
By storing all session state in PostgreSQL JSONB rather than
application memory, the system supports session resumption after
network interruption---a critical property for rural deployment where
connectivity may be intermittent. This also enables horizontal scaling
without sticky session requirements.

\paragraph{Dual-path transcription.}
Providing both browser-native Web Speech API and server-side
\texttt{faster-whisper} allows the system to adapt to connectivity
conditions. The browser path requires no server round-trip, while the
server path provides better accuracy for low-microphone-quality audio
when network conditions allow upload.

\paragraph{Confirmation gate calibration.}
The dual-threshold gate in \eqref{eq:confirm} is motivated by the
observation, supported by voice-based agricultural advisory
research~\cite{ref3,ref20}, that excessive confirmation dialogs are a
primary driver of voice-interface abandonment. The design suppresses
confirmations on clear inputs while maintaining accuracy on
ambiguous ones.

\paragraph{Scenario monotonicity.}
The verified property that pessimistic $<$ base $<$ optimistic holds
for all revenue-derived metrics is a critical reliability guarantee
for operators. If the risk envelope communicated by the system is
internally inconsistent, operators cannot trust it for investment
decisions.

\paragraph{Performance design targets.}
Table~\ref{tab:perf} lists the response-time acceptance criteria that
will be used as pass/fail thresholds in the planned
$N{\geq}30$-session load test during the field-trial phase. These
targets have not yet been measured under load and represent engineering
design goals informed by agricultural mobile-application usability
guidelines~\cite{ref20}.

\begin{table}[tbp]
  \caption{Performance Design Targets for Field-Trial Load Testing}
  \label{tab:perf}
  \centering
  \small
  \begin{tabular}{lcc}
    \toprule
    \textbf{Component} & \textbf{Metric} & \textbf{Target} \\
    \midrule
    Session answer endpoint   & P95   & $<$200\,ms \\
    Server-side transcription & P95   & $<$5\,s    \\
    Financial analysis        & P95   & $<$2\,s    \\
    Market price fetch        & P95   & $<$3\,s    \\
    E2E survey to report      & Total & $<$3\,min  \\
    \bottomrule
    \multicolumn{3}{l}{\footnotesize
      P95: 95th-percentile response time.}\\
  \end{tabular}
\end{table}

\noindent\textit{Limitation---synthetic validation.}
All current tests use constructed fixtures. No user study has yet been
conducted. SUS score, real-world word-error rate, and projection
accuracy against realised farm outcomes remain unmeasured and are
priorities for the field-trial phase.

\medskip
\noindent\textit{Limitation---English-only STT.}
The \texttt{faster-whisper} pipeline has not been benchmarked on
Indian English accent diversity or regional language input.
Multilingual support is a priority for production deployment,
especially given the pace of digital agriculture adoption in
India~\cite{ref18}.

\subsection{Conclusion}

This paper presents AquaponicAI, a voice-driven decision-support
platform for which a 36-case validation suite was executed with
full pass rate on all deterministic assertions in our evaluation run.
The platform is designed to enable farm operators to obtain
personalised three-scenario financial plans via voice or text input,
without requiring financial expertise; end-to-end completion time
is targeted at under three minutes, pending field-trial verification.
The modular, containerized architecture is intended to be deployable
in resource-constrained agricultural contexts.

Four directions are identified for future work:
(i)~field trials with $N{\geq}50$ operators measuring task completion
rate, SUS score, and projected-vs-realised ROI over 12 months;
(ii)~benchmarking the ASR pipeline against IndicWav2Vec for accent
robustness;
(iii)~replacing the rule-based recommendation engine with a gradient
boosting model trained on historical farm financial outcomes; and
(iv)~extending real-time market price integration to commodity futures
for forward-looking scenario bounds.

% ── ACKNOWLEDGMENT & DISCLOSURE ─────────────────────────────────────────────
\begin{credits}
\subsubsection{\ackname}
The authors thank the Department of MCA, RV College of Engineering,
Bengaluru, for providing the research infrastructure that supported
this work.

\subsubsection{\discintname}
The authors have no competing interests to declare that are relevant
to the content of this article.
\end{credits}


% ── REFERENCES ──────────────────────────────────────────────────────────────
\begin{thebibliography}{28}

\bibitem{ref1}
Verma, A.K., Chandrakant, M.H., John, V.C., Peter, R.M., John, I.E.:
Aquaponics as an integrated agri-aquaculture system (IAAS): Emerging
trends and future prospects.
Technol.\ Forecast.\ Soc.\ Chang.\ \textbf{194}, 122709 (2023).
\doi{10.1016/j.techfore.2023.122709}

\bibitem{ref2}
Durai, S.K.S., Shamili, M.D.:
Smart farming using machine learning and deep learning techniques.
Decis.\ Anal.\ J.\ \textbf{3}, 100041 (2022).
\doi{10.1016/j.dajour.2022.100041}

\bibitem{ref3}
Biswas, R., Goel, N.:
Intelligent chatbot assistant in agriculture domain.
In: Proc.\ Int.\ Conf.\ Agriculture-Centric Computation,
pp.\ 180--194. Springer, Cham (2023)

\bibitem{ref4}
Lobanov, V., Pate, J., Latrille, F., Joyce, A.:
Aquaponics and the crossroads of profitability and sustainability
viewed through sensitivity and risk modeling.
J.\ Clean.\ Prod.\ \textbf{528}, 146674 (2025).
\doi{10.1016/j.jclepro.2025.146674}

\bibitem{ref5}
Radford, A., Kim, J.W., Xu, T., Brockman, G., McLeavey, C.,
Sutskever, I.:
Robust speech recognition via large-scale weak supervision.
In: Proc.\ 40th Int.\ Conf.\ Machine Learning (ICML), vol.\ 202,
pp.\ 28492--28518. PMLR (2023)

\bibitem{ref7}
Okomoda, V.T., Oladimeji, S.A., Solomon, S.G., Olufeagba, S.O.,
Ogah, S.I., Ikhwanuddin, M.:
Aquaponics production system: A review of historical perspective,
opportunities, and challenges of its adoption.
Food Sci.\ Nutr.\ \textbf{11}(3), 1157--1165 (2023).
\doi{10.1002/fsn3.3154}

\bibitem{ref9}
Veena, G., Chandana, S., Nandini, A., Mayuri, S., Kiran, D.R.,
Sowmya, K.:
AGRONER: Agricultural named entity recognition using pre-trained
language models.
Expert Syst.\ Appl.\ \textbf{232}, 120440 (2023).
\doi{10.1016/j.eswa.2023.120440}

\bibitem{ref10}
Baevski, A., Zhou, Y., Mohamed, A., Auli, M.:
wav2vec 2.0: A framework for self-supervised learning of speech
representations.
In: Proc.\ Adv.\ Neural Inf.\ Process.\ Syst.\ (NeurIPS),
vol.\ 33, pp.\ 12449--12460 (2020)

\bibitem{ref11}
Mohanty, S.P., Rout, B.K., Satpathy, B.:
Forecasting agricultural commodity prices using a novel model based
on NARX and SVR.
Neural Comput.\ Appl.\ \textbf{35}(12), 8661--8676 (2023).
\doi{10.1007/s00521-023-08528-7}

\bibitem{ref15}
Yousaf, A., Mehmood, A., Sang-Ary, G., Mehmood, P., Kim, H.S.:
Precision agriculture utilizing blockchain, edge, and fog computing:
A systematic review of the state of the art.
Front.\ Sustain.\ Food Syst.\ \textbf{6}, 1053921 (2022).
\doi{10.3389/fsufs.2022.1053921}

\bibitem{ref17}
Paul, A., Ghosh, D., Bhattacharjee, A., Chattopadhyay, S.,
Majumdar, S.:
Forecasting agricultural commodity prices: A machine learning approach
on Indian market data.
PLOS ONE \textbf{17}(7), e0270553 (2022).
\doi{10.1371/journal.pone.0270553}

\bibitem{ref18}
Balkrishna, A.A., Kaushik, V.K., Kumar, K., Bhatt, M., Misra, A.:
Significance of smart agriculture technologies in India: Initiatives,
challenges, and future directions.
Smart Agric.\ Technol.\ \textbf{5}, 100318 (2023).
\doi{10.1016/j.atech.2023.100318}

\bibitem{ref20}
Deshmukh, A., Chalmeta, R.:
User experience and usability of voice user interfaces: A systematic
literature review.
Information \textbf{15}(9), 579 (2024).
\doi{10.3390/info15090579}

\bibitem{ref21}
Goddek, S., Joyce, A., Kotzen, B., Burnell, G.M. (eds.):
Aquaponics Food Production Systems: Combined Aquaculture and
Hydroponic Production Technologies for the Future.
Springer, Cham (2019). \doi{10.1007/978-3-030-15943-6}

\bibitem{ref22}
Somerville, C., Cohen, M., Pantanella, E., Stankus, A., Lovatelli, A.:
Small-scale aquaponic food production: Integrated fish and plant farming.
FAO Fisheries and Aquaculture Technical Paper No.\ 589. FAO, Rome (2014)

\bibitem{ref23}
Kamilaris, A., Prenafeta-Bol\'{u}d, F.X.:
Deep learning in agriculture: A survey.
Comput.\ Electron.\ Agric.\ \textbf{147}, 70--90 (2018).
\doi{10.1016/j.compag.2018.02.016}

\bibitem{ref24}
Wolfert, S., Ge, L., Verdouw, C., Bogaardt, M.J.:
Big data in smart farming --- a review.
Agric.\ Syst.\ \textbf{153}, 69--80 (2017).
\doi{10.1016/j.agsy.2017.01.023}

\bibitem{ref25}
Devlin, J., Chang, M.W., Lee, K., Toutanova, K.:
BERT: Pre-training of deep bidirectional transformers for language
understanding. In: Proc.\ 2019 Conf.\ North American Chapter of the
Association for Computational Linguistics (NAACL-HLT), pp.\ 4171--4186.
ACL, Minneapolis (2019)

\bibitem{ref26}
Graves, A., Mohamed, A.R., Hinton, G.:
Speech recognition with deep recurrent neural networks.
In: Proc.\ IEEE Int.\ Conf.\ Acoustics, Speech and Signal Processing
(ICASSP), pp.\ 6645--6649. IEEE (2013).
\doi{10.1109/ICASSP.2013.6638947}

\bibitem{ref27}
Liakos, K.G., Busato, P., Moshou, D., Pearson, S., Bochtis, D.:
Machine learning in agriculture: A review.
Sensors \textbf{18}(8), 2674 (2018). \doi{10.3390/s18082674}

\bibitem{ref28}
Park, D.S., Chan, W., Zhang, Y., Chiu, C.C., Zoph, B.,
Cubuk, E.D., Le, Q.V.:
SpecAugment: A simple data augmentation method for automatic speech
recognition. In: Proc.\ Interspeech, pp.\ 2613--2617 (2019).
\doi{10.21437/Interspeech.2019-2680}

\bibitem{ref29}
Vaswani, A., Shazeer, N., Parmar, N., Uszkoreit, J., Jones, L.,
Gomez, A.N., Kaiser, {\L}., Polosukhin, I.:
Attention is all you need.
In: Proc.\ Adv.\ Neural Inf.\ Process.\ Syst.\ (NeurIPS),
vol.\ 30, pp.\ 5998--6008 (2017)

\bibitem{ref30}
Adamopoulou, E., Moussiades, L.:
An overview of chatbot technology.
In: Proc.\ IFIP Int.\ Conf.\ Artificial Intelligence Applications
and Innovations (AIAI), pp.\ 373--383. Springer, Cham (2020).
\doi{10.1007/978-3-030-49186-4\_31}

\bibitem{ref31}
Quagrainie, K.K., Flores, R.M.V., Kim, H.J., McClain, V.:
Economic analysis of aquaponics and hydroponics production in
the US Midwest.
J.\ Appl.\ Aquac.\ \textbf{30}(1), 1--14 (2018).
\doi{10.1080/10454438.2017.1414009}

\bibitem{ref32}
Saiz-Rubio, V., Rovira-M\'{a}s, F.:
From smart farming towards agriculture 5.0: A perspective from
mastery thinking.
Agronomy \textbf{10}(2), 157 (2020). \doi{10.3390/agronomy10020157}

\bibitem{ref33}
Brealey, R.A., Myers, S.C., Allen, F.:
Principles of Corporate Finance. 13th edn.
McGraw-Hill Education, New York (2020)

\bibitem{ref34}
Bobadilla, J., Ortega, F., Hernando, A., Guti\'{e}rrez, A.:
Recommender systems survey.
Knowl.-Based Syst.\ \textbf{46}, 109--132 (2013).
\doi{10.1016/j.knosys.2013.03.012}

\bibitem{ref35}
Birthal, P.S., Roy, D., Negi, D.S.:
Assessing the impact of crop diversification on farm income of
smallholder farmers: A case study from India.
World Dev.\ \textbf{75}, 28--45 (2015).
\doi{10.1016/j.worlddev.2015.01.015}

\end{thebibliography}

\end{document}
